\documentclass[11pt]{article}
\usepackage[T1]{fontenc}
\usepackage{lmodern}
\usepackage[margin=1in]{geometry}
\usepackage{amsmath,amssymb}
\usepackage{microtype}
\usepackage{booktabs}
\usepackage{xcolor}
\usepackage[colorlinks=true,linkcolor=blue!50!black]{hyperref}

\newcommand{\C}{\mathbb{C}}
\newcommand{\abs}[1]{\lvert #1\rvert}
\newcommand{\norm}[1]{\lVert #1\rVert}
\newcommand{\key}[1]{\par\smallskip\noindent{\bfseries #1}\par\nobreak\vspace{-2pt}}
\setlength{\abovedisplayskip}{4pt}\setlength{\belowdisplayskip}{4pt}
\setlength{\abovedisplayshortskip}{2pt}\setlength{\belowdisplayshortskip}{2pt}

\title{\bfseries Step 4 on one example\\[4pt]
\large\normalfont the two graphs on the candidates of the Fourier pair in $\C^6$, and the pattern that is not there}
\author{}\date{}

\begin{document}
\maketitle
\vspace{-2.6em}

\noindent
Step 4 decides the last question with a graph. Inside one certified region the first two bases are
known only approximately; a vector unbiased to both is a \emph{candidate}; candidates are grouped into
balls; and two graphs are built on the balls, with $O$ joining those that might be orthogonal and $U$
those that might be unbiased. A third and a fourth basis would appear as two disjoint six-cliques in
$O$ with all $36$ cross pairs in $U$. Because $O$ and $U$ are \emph{supersets} of the true relations,
absence of that pattern is a proof and presence of it is not.

\smallskip
\noindent
The example runs the whole of that on the most symmetric pair of bases in $\C^6$ --- the standard basis
and the Fourier basis $F/\sqrt6$ --- because there the candidate set can be computed exactly rather
than sampled.

\key{1.\ The candidates, enumerated exactly}
Unbiasedness to the standard basis forces $\abs{x_k}=1/\sqrt6$; write $x_k=\zeta^{a_k}/\sqrt6$ with
$\zeta=e^{i\pi/6}$. Unbiasedness to the $j$-th Fourier column says
\[
\abs{\langle f_j,x\rangle}=\frac16\abs{\textstyle\sum_k\zeta^{a_k}\omega^{-jk}}=\frac1{\sqrt6}
\quad\Longleftrightarrow\quad
\Bigl|\textstyle\sum_k\zeta^{a_k-2jk}\Bigr|^2=6 ,
\]
and a sum of twelfth roots of unity is compared with $6$ exactly, since
$\abs{\sum_c n_c\zeta^c}^2=A+B\sqrt3$ with $A,B$ rational and computable from the counts $n_c$.
Fixing $a_0=0$ (global phase) and enumerating all $12^5=248{,}832$ exponent tuples gives
\[
\textbf{12 candidate rays},\qquad 144\ \text{vectors counting all global phases},
\]
with exponent tuples $a_0,\dots,a_5$
\begin{center}\small
\begin{tabular}{@{}lll@{}}
$(0,1,0,9,4,9)$ & $(0,1,4,9,4,1)$ & $(0,3,4,3,0,7)$\\
$(0,3,8,3,0,11)$ & $(0,5,0,9,8,9)$ & $(0,5,8,9,8,5)$\\
$(0,7,0,3,4,3)$ & $(0,7,4,3,4,7)$ & $(0,9,4,9,0,1)$\\
$(0,9,8,9,0,5)$ & $(0,11,0,3,8,3)$ & $(0,11,8,3,8,11)$
\end{tabular}
\end{center}
Every one of the $72$ checks $\abs{\cdot}^2=6$ was done in exact arithmetic; the residual is the
integer $0$.

\emph{A detail worth knowing.} On the \emph{sixth}-root lattice $\zeta^{a}=\omega^{a/2}$ the same
enumeration returns \textbf{no candidate at all}, because there $\abs{\sum}^2$ is a value of the norm
form $p^2+pq+q^2$ and $6$ is not one. The candidates appear only when half-integer powers of
$\omega$ are allowed. A search that assumes sixth-root phases would conclude, wrongly, that this pair
of bases admits no third basis.

\key{2.\ The graph $O$: 48 edges, and four bases}
Two candidates are joined in $O$ when the inner product is \emph{exactly} zero (checked in
$\mathbb{Z}[\zeta]$, by reducing $\sum_k\zeta^{a_k-b_k}$ to the basis $1,\zeta,\zeta^2,\zeta^3$ and
testing four integer coefficients). The result is a very regular graph:
\begin{center}\small
\begin{tabular}{@{}lr@{}}
\toprule
candidates & $12$\\
pairs of candidates & $66$\\
edges of $O$ (exactly orthogonal pairs) & $48$\\
degree of every vertex & $8$ \quad(each candidate is orthogonal to $8$ of the other $11$)\\
six-cliques in $O$ (orthonormal bases) & $\mathbf{4}$\\
seven-cliques in $O$ & $0$\\[2pt]
\multicolumn{2}{@{}l}{\small the four bases, as candidate indices:}\\
& $\{0,2,5,6,8,11\},\ \ \{0,3,5,7,8,10\},\ \ \{1,2,4,6,9,11\},\ \ \{1,3,4,7,9,10\}$\\
\bottomrule
\end{tabular}
\end{center}
Four bases, and only four: since $4\times6=24=2\times12$, each candidate belongs to exactly two of
them. This is the structure the doc warns about --- single candidates are plentiful and bases are not.

\key{3.\ The graph $U$: empty, and the pattern with it}
Two candidates are joined in $U$ when $\abs{\langle x,y\rangle}=1/\sqrt6$ exactly. The answer is
striking:
\[
\textbf{0 edges in }U\text{ among the }66\text{ pairs},
\qquad\text{and the }66\text{ inner products take exactly two moduli:}
\]
$\abs{\langle x,y\rangle}=0$ for $48$ pairs and $\abs{\langle x,y\rangle}=1/\sqrt3$ for the other
$18$. There is no third value and nothing in between: the candidates of this pair sit at one of two
distances from each other, and $1/\sqrt3$ is as close to $1/\sqrt6$ as any pair ever gets.
So the fourth basis cannot be assembled, and it fails at the first opportunity: not merely do the four
bases fail to be pairwise unbiased, no two candidates whatever are unbiased. Of the six pairs of
bases, two are disjoint as sets of vectors, and each of those has
\[
0\ \text{of}\ 36\ \text{cross pairs unbiased}
\]
--- the pattern of Step~4 is absent from this candidate set, and the region is excluded.

The distance to the threshold is worth putting in the run's own units. The edge test forms
$t=\abs{\langle x,y\rangle}$ and compares
\[
\text{disto}=\frac{\sqrt6\,t}{2\pi}\quad\text{and}\quad
\text{distu}=\frac{\abs{\sqrt6\,t-1}}{2\pi}
\]
with the ball radius. For a pair of candidates taken from one basis, $t=2.9\times10^{-17}$, so
$\text{disto}=1.1\times10^{-17}$: the $O$ edge is kept with the whole margin $2\times10^{-7}$. For the
\emph{closest} of the $36$ cross pairs, $t=1/\sqrt3=0.57735$, so
\[
\text{distu}=0.065924 ,
\]
which is $3\times10^{5}$ times the largest threshold any ball of the run's size could give. The
unbiasedness test is not failing marginally here; it is failing by two thirds of the range.

\key{4.\ What a radius does to an edge}
The graphs are built on balls, not points, and an edge is kept unless it can be \emph{refuted}. Rotate
one entry of a vector of a basis by $\delta$ cycles: the centre moves by
$\sqrt2\,\abs{e^{2\pi i\delta}-1}/\sqrt6$, and the $O$ edge to its orthogonal partner decays linearly
in $\delta$:
\begin{center}\small
\begin{tabular}{@{}r r r l r@{}}
\toprule
$\delta$ (cycles) & displacement of the centre & $\text{disto}$ & $O$ edge at radius $0$ & radius that keeps it\\
\midrule
$10^{-7}$ & $2.57\times10^{-7}$ & $4.08\times10^{-8}$ & kept & $0$\\
$10^{-6}$ & $2.57\times10^{-6}$ & $4.08\times10^{-7}$ & \textbf{refused} & $2.08\times10^{-7}$\\
$10^{-5}$ & $2.57\times10^{-5}$ & $4.08\times10^{-6}$ & refused & $3.88\times10^{-6}$\\
$10^{-4}$ & $2.57\times10^{-4}$ & $4.08\times10^{-5}$ & refused & $4.06\times10^{-5}$\\
\bottomrule
\end{tabular}
\end{center}
Two readings. The exclusion is sharp: a single entry rotated by one part in $10^6$ of a turn already
destroys the orthogonality, so at radius zero the graph is exactly the truth. And the exclusion is
fragile in the right direction: with balls of radius $2\times10^{-7}$ the same edge is kept, because
the ball might contain an exactly orthogonal pair. That is a false positive --- permitted, since $O$
must be a superset --- and it is why the answer at a coarse region is always ``unresolved'', never
``excluded''.

\key{5.\ The other thirty-six candidates}
The enumeration of \S1 is exhaustive only over twelfth roots of unity, and nothing forces a candidate
onto a lattice: the condition is two modulus conditions on six continuous phases. In fact the lattice
set is a proper subset. A multistart search of the six phase equations ($4000$ random starts, final
residual below $10^{-14}$; \texttt{step4\_example.py -{}-slow}) finds
\[
\textbf{48}\ \text{candidates up to global phase},\qquad
300\ \text{exactly orthogonal pairs},\qquad 16\ \text{orthonormal bases},
\]
and still \textbf{no} pair of candidates unbiased: the pattern is absent there too. This is the count
quoted in \texttt{mub-step4.pdf} from the literature, and the two enumerations agree once the lattice
is dropped: the search puts exactly $12$ of the $48$ on the twelfth-root lattice --- the same twelve
that \S1 enumerates --- and the other $36$ on none of the lattices of order $6,12,24$ or $48$, which
is why no lattice enumeration could have found them.

The consequence is worth stating plainly, because it cuts both ways. The graph of \S2--3 is a genuine
\emph{subgraph} of the real one: its four bases are real bases and its empty $U$ is really empty on
those twelve vectors, so the exclusion it supports is sound --- but a lattice-restricted search is
blind to thirty-six vertices, and a pattern living on them would not be seen. Coverage, not direction,
is what a lattice costs; that is why the run's Stage~C works with balls and real iteration rather than
with a lattice.

\key{6.\ What these numbers do not say}
Everything above is for \emph{one} pair of bases and for candidates on the twelfth-root lattice. The
enumeration is exhaustive \emph{on that lattice} and therefore cannot rule out candidates with
arbitrary phases; the count $12$ is a lower bound on the candidate set, and a larger set would only add
vertices and edges, which is harmless for the direction of the implication but would change every
number in the tables. In the run, Stage~C does this for $23{,}506$ regions with balls, not points, and
its graph has hundreds of vertices, not twelve or forty-eight. Finally, the absence of the pattern here is a fact
about this candidate list; it becomes a proof only when the list is known to contain every candidate
of the region, which is what Step~3's certified balls provide.

\medskip
\noindent\footnotesize\textit{Provenance.} The two graphs, the pattern, and the edge tests are Step~4's
as described in \texttt{mub-step4.pdf} and preprint \S6; the pairtest constants ($10^{-10}$, $2\times
10^{-7}$, the factor $\sqrt6$) are those of the documented implementation \texttt{stageC.cpp}. Every
number above was printed by \texttt{examples/step4\_example.py}, which carries its own exact
arithmetic in $\mathbb{Z}[\zeta_{12}]$ and takes about one second; the $48$-candidate count of \S5
comes from a multistart Newton search, \texttt{examples/step4\_example.py -{}-slow}, and reproduces
the literature count quoted in \texttt{mub-step4.pdf}.

\end{document}
